% Standalone problem document, generated from the adjacent README.md.
% Compile with XeLaTeX. Mathematical content is maintained in Markdown.
\documentclass[11pt,a4paper]{article}
\usepackage[fontsize=10.5pt]{scrextend}
\usepackage[margin=22mm,headheight=15pt,headsep=9mm,footskip=11mm]{geometry}
\usepackage{fontspec}
\setmainfont{texgyrepagella}[Extension=.otf,UprightFont=*-regular,BoldFont=*-bold,ItalicFont=*-italic,BoldItalicFont=*-bolditalic]
\setsansfont{texgyreheros}[Extension=.otf,UprightFont=*-regular,BoldFont=*-bold,ItalicFont=*-italic,BoldItalicFont=*-bolditalic]
\setmonofont{texgyrecursor}[Extension=.otf,UprightFont=*-regular,BoldFont=*-bold,ItalicFont=*-italic,BoldItalicFont=*-bolditalic,Scale=MatchLowercase]
\usepackage{amsmath,amssymb}
\usepackage{unicode-math}
\setmathfont{texgyrepagella-math.otf}
\usepackage{xcolor,microtype,enumitem,needspace,fancyhdr}
\usepackage{bookmark}
\definecolor{ink}{HTML}{17344B}
\definecolor{accent}{HTML}{197D80}
\definecolor{muted}{HTML}{596773}
\hypersetup{colorlinks=true,linkcolor=accent,urlcolor=accent,pdftitle={MF-04 - Finiteness
for nonnegative rational matrix
families},pdfauthor={Open Problems in Numerical Linear Algebra}}
\urlstyle{same}
\setlength{\parindent}{0pt}
\setlength{\parskip}{5pt plus 1pt minus 1pt}
\linespread{1.04}
\setlength{\emergencystretch}{3em}
\widowpenalty=10000
\clubpenalty=10000
\displaywidowpenalty=10000
\allowdisplaybreaks[1]
\setlist{nosep,leftmargin=1.5em}
\providecommand{\tightlist}{\setlength{\itemsep}{0pt}\setlength{\parskip}{0pt}}
\providecommand{\pandocbounded}[1]{#1}
\pagestyle{fancy}
\fancyhf{}
\fancyhead[L]{\sffamily\footnotesize\color{muted}OPEN PROBLEMS IN NUMERICAL LINEAR ALGEBRA}
\fancyhead[R]{\sffamily\bfseries\footnotesize\color{accent}MF-04}
\fancyfoot[L]{\parbox[b]{0.9\textwidth}{\sffamily\fontsize{8}{10}\selectfont\color{muted}Editorial ratings; a bounded search cannot certify that a problem remains unsolved.\\Literature check: 2026-09-10}}
\fancyfoot[R]{\sffamily\footnotesize\color{muted}\thepage}
\renewcommand{\headrulewidth}{0.3pt}
\renewcommand{\headrule}{\hbox to\headwidth{\color{accent}\leaders\hrule height \headrulewidth\hfill}}
\makeatletter
\renewcommand{\section}{\Needspace{5\baselineskip}\@startsection{section}{1}{0pt}{12pt plus 2pt minus 2pt}{4pt}{\sffamily\bfseries\large\color{ink}}}
\renewcommand{\subsection}{\Needspace{4\baselineskip}\@startsection{subsection}{2}{0pt}{9pt plus 1pt minus 1pt}{3pt}{\sffamily\bfseries\normalsize\color{ink}}}
\makeatother
\setcounter{secnumdepth}{0}
\begin{document}
{\sffamily\small\color{muted}Matrix functions and stability\par}
\vspace{3pt}
{\raggedright\hyphenpenalty=10000\exhyphenpenalty=10000\sffamily\bfseries\fontsize{21}{24}\selectfont\color{ink}Finiteness
for nonnegative rational matrix families\par}
\vspace{7pt}
{\sffamily\small\textbf{Difficulty:} extreme\quad\textbf{Importance:} broadly
interesting\par}
{\sffamily\small\bfseries\color{accent}Status: Partially resolved\par}
\vspace{4pt}
{\color{accent}\hrule height 0.8pt}
\vspace{6pt}
\textbf{Rating rationale:} Extreme because rational-input finiteness is
a longstanding barrier in joint spectral-radius theory; its consequences
reach control, combinatorics, and exact computation.

\subsection{Context and notation}\label{context-and-notation}

The joint spectral radius of a nonempty compact set
\(\mathcal M\subset\mathbb C^{d\times d}\) is

\[
\widehat\rho(\mathcal M)=\lim_{k\to\infty}
\max_{A_1,\ldots,A_k\in\mathcal M}\|A_k\cdots A_1\|_2^{1/k}.
\]

This definition also applies to finite real matrix sets. The ordinary
spectral radius of one matrix is written \(\rho(A)\).

\subsection{Problem statement}\label{problem-statement}

Does every finite nonempty set
\(\mathcal M\subset\mathbb Q_{\ge0}^{d\times d}\), for every positive
integer \(d\), admit a positive integer \(k\) and matrices
\(A_1,\ldots,A_k\in\mathcal M\) such that

\[
\widehat\rho(\mathcal M)=\rho(A_k\cdots A_1)^{1/k}?
\]

The entries must be nonnegative rationals; unrestricted real entries
change the status. This asks whether a finite product realizes the
asymptotic growth rate.

\subsection{References and status
evidence}\label{references-and-status-evidence}

Jungers and Blondel, \href{https://arxiv.org/abs/math/0702489}{On the
finiteness property for rational matrices}, Linear Algebra Appl. 428
(2008), 2283--2295,
\href{https://doi.org/10.1016/j.laa.2007.07.007}{DOI}, abstract and
reduction theorems; Jungers, \emph{The Joint Spectral Radius: Theory and
Applications},
\href{https://doi.org/10.1007/978-3-540-95980-9_5}{Chapter 4},
pp.~63--74 (2009). They reduce this question to pairs of binary matrices
in arbitrary dimension. Mejstrik,
\href{https://arxiv.org/pdf/2505.10178}{The finiteness conjecture for
3×3 binary matrices}, §3 and Theorem 3.1, \emph{Dolomites Research Notes
on Approximation} 15 (2022), 24--38 (deposited on arXiv in 2025), treats
the all-dimensions problem as open while proving the binary-pair case in
dimension three. Searches for rational finiteness and binary-pair
counterexamples through the check date found no resolution of the
general question. Equivalent binary formulations count once.

\subsection{Audit --- 2026-09-10}\label{audit-2026-09-10}

Rechecked the
\href{https://drna.padovauniversitypress.it/system/files/papers/MEJSTRIK.pdf}{published
Mejstrik paper}, §3: binary pairs through dimension three are settled,
while arbitrary dimensions remain open. Corrected its publication date
to 2022; the \href{https://arxiv.org/abs/2505.10178}{2025 arXiv date} is
a deposit date. Rational-finiteness and later binary-pair searches found
no general resolution.
\end{document}
